\begin{afterword}
\summaryhead

The post starts from the slowness of neurons. Because a neuron fires at most a few hundred times a second, the brain can take only about a hundred sequential steps in any quick mental operation. The author infers that the brain must rely heavily on caching: storing earlier results and looking them up instead of recomputing them. Most human thought, the post guesses, consists of such lookups.

The post then applies the idea to beliefs taken from other people. A man once knocked out his chimney from the bottom up because a neighbor had said that was the way to do it. The author suspects the man would have examined the idea if it had been his own. Civilization depends on such borrowing, the essay grants, but people who think of themselves as critical repeat slogans they never examined: that religion cannot be disproved, that death gives life meaning, that love is not rational. Readers are asked to stop their minds from completing such patterns and to think.

\responsehead

The post gives its readers a label and no method. The label is ``cached thought.'' The method would be a way to tell which borrowed conclusions deserve rechecking, and the post concedes that most must be borrowed. It never says which.

The post changes topic partway through. The first four paragraphs are about neurons, spike rates and a rule from cognitive science. They make the idea sound like a finding about the brain. The rest of the post is about something else: people believing what they were told. Trust in testimony has nothing to do with how fast neurons fire, and nothing in the opening bears on it. The word ``cache'' makes the second topic sound as scientific as the first.

Its one psychological claim is contradicted by the evidence. The post assumes that people scrutinize their own ideas more than borrowed ones. An experiment on how people evaluate arguments found the opposite, and the author's own posts of the same fortnight, on rationalization and the bottom line, are built on the opposite. The main example is a story the author cannot source.

The post's examples show how the label is used. Every example of a cached thought is an opinion the author rejects: that religion is beyond evidence, that death gives meaning, that the species may not deserve to survive, that love is not rational. One of them stands in the way of the author's own cause, which the post names: existential risk. The people who hold these views are ``good and decent folk'' in whom ``the brain completes the pattern.'' The reader is the one who thinks. The final paragraph says who the label is for: the next time someone repeats ``a meme you think is silly or false.'' The reader has already judged the view silly or false before the label is applied, so the label adds no reason. It confirms a verdict already reached and supplies a way to dismiss the person who holds the view.

In short: a neuroscience frame that supports nothing in the argument, a psychological claim that the evidence contradicts, and a label that lets readers dismiss views they already reject without saying why.
\end{afterword}
